%
% ausgangskreis.tex -- Ausgangskreis
%
% (c) 2021 Prof Dr Andreas Müller, OST Ostschweizer Fachhochschule
%
\documentclass[tikz]{standalone}
\usepackage{amsmath}
\usepackage{times}
\usepackage{txfonts}
\usepackage{pgfplots}
\usepackage{csvsimple}
\usetikzlibrary{arrows,intersections,math}
\begin{document}
\def\skala{1}
\begin{tikzpicture}[>=latex,thick,scale=\skala]

% Kreis k0: Mittelpunkt M0, Radius 1.8
\def\R{1.8}
\pgfmathparse{1.4*\R}
\xdef\R{\pgfmathresult}
\coordinate (M0) at (0, 0);
\coordinate (M1) at (0, \R);
\coordinate (Z)  at (0,-\R);
% Kreis
\draw (M0) circle (\R);
% Vertikale Linie M1--Z
\draw (M1) -- (Z);
% Punkte – Labels mit weissem Hintergrund damit Linie nicht durchscheint
\fill (M0) circle (2pt) node[right=6pt, fill=white, inner sep=1pt] {$M_0$};
\fill (M1) circle (2pt) node[right=6pt, fill=white, inner sep=1pt] {$M_1$};
\fill (Z)  circle (2pt) node[right=6pt, fill=white, inner sep=1pt] {$Z$};

\node[fill=white, inner sep=1pt] at (-1.86, 2.14) {$k_0$};

\end{tikzpicture}
\end{document}

